% Generated from lessons/0009-3-positive-series.html; edit the HTML source, then rerun the converter.
\section{正项级数}
\lessonmark{正项级数}

这一节是数项级数判别法的核心。只要每一项非负，部分和就单调递增，于是“收敛”变成了“能不能被上界压住”。后面的比较判别、根值判别、比值判别、Raabe 判别和积分判别，本质上都是在找一个能压住它的参照物，或证明它压不住。

\subsection*{本节主线}

\begin{conceptbox}
\textbf{先看正项}
若 \(x_n\ge 0\)，部分和 \(S_n=\sum_{k=1}^n x_k\) 单调递增。
\end{conceptbox}

\begin{conceptbox}
\textbf{再找参照}
和 \(p\)-级数、几何级数、阶乘型、指数型或反常积分比较。
\end{conceptbox}

\begin{conceptbox}
\textbf{最后选判别法}
幂指数用根值，阶乘递推用比值，比值临界再试 Raabe 或积分判别。
\end{conceptbox}

\begin{notebox}
考试顺序：先判是否正项；再看一般项是否趋于 \(0\)；然后按“比较、根值、比值、Raabe、积分”逐层尝试。
\end{notebox}

\subsection*{一、定义与收敛原理}

\subsubsection*{定义 9.3.1：正项级数}

若级数 \(\sum_{n=1}^{\infty}x_n\) 的每一项都非负：

\[
      x_n\ge 0,\qquad n=1,2,\ldots,
    \]

则称它为正项级数。

\subsubsection*{定理 9.3.1：正项级数的收敛原理}

正项级数 \(\sum x_n\) 收敛，当且仅当它的部分和数列

\[
      S_n=x_1+x_2+\cdots+x_n
    \]

有上界。若 \(\{S_n\}\) 无上界，则 \(\sum x_n\) 发散到 \(+\infty\)。

\paragraph*{证明思路}

因为 \(x_{n+1}\ge 0\)，所以 \(S_{n+1}=S_n+x_{n+1}\ge S_n\)。于是 \(\{S_n\}\) 是单调递增数列。单调递增数列收敛当且仅当有上界；没有上界就趋于 \(+\infty\)。

\begin{notebox}
正项级数的本质：不可能靠“正负抵消”收敛，只能靠每一项足够快地变小，让部分和被上界压住。
\end{notebox}

\subsection*{二、比较判别法}

\subsubsection*{定理 9.3.2：比较判别法}

设 \(\sum x_n,\sum y_n\) 都是正项级数。若存在常数 \(A>0\)，使得从某一项起有

\[
      0\le x_n\le A y_n,
    \]

则：

\begin{itemize}
 \item 若 \(\sum y_n\) 收敛，则 \(\sum x_n\) 收敛；
 \item 若 \(\sum x_n\) 发散，则 \(\sum y_n\) 发散。
 \end{itemize}

\paragraph*{证明思路}

设 \(S_n=\sum_{k=1}^n x_k,\ T_n=\sum_{k=1}^n y_k\)。由 \(x_n\le A y_n\) 得 \(S_n\le A T_n+\text{有限常数}\)。如果 \(T_n\) 有界，则 \(S_n\) 有界，正项级数收敛。

\begin{notebox}
有限多项不影响级数敛散性，所以比较不一定要从第一项开始，只要“最终成立”即可。
\end{notebox}

\subsubsection*{定理 9.3.2'：极限比较判别法}

设 \(\sum x_n,\sum y_n\) 是正项级数，且

\[
      \lim_{n\to\infty}\frac{x_n}{y_n}=l,\qquad 0\le l\le+\infty.
    \]

\begin{itemize}
 \item 若 \(0\le l<+\infty\)，且 \(\sum y_n\) 收敛，则 \(\sum x_n\) 收敛；
 \item 若 \(0<l\le+\infty\)，且 \(\sum y_n\) 发散，则 \(\sum x_n\) 发散；
 \item 若 \(0<l<+\infty\)，则 \(\sum x_n\) 与 \(\sum y_n\) 同敛散。
 \end{itemize}

\paragraph*{证明思路}

若 \(l<+\infty\)，则充分大的 \(n\) 有 \(x_n/y_n<l+1\)，所以 \(x_n<(l+1)y_n\)。若 \(l>0\)，则充分大的 \(n\) 有 \(x_n/y_n>l/2\)，所以 \(x_n>(l/2)y_n\)。剩下交给比较判别法。

\subsubsection*{典型例子}

\begin{conceptbox}
\textbf{有理式}
\[
          \sum_{n=1}^{\infty}\frac{n+3}{2n^3-n}
        \]

与 \(\sum 1/n^2\) 比较，因为一般项约为 \(\frac{1}{2n^2}\)，所以收敛。
\end{conceptbox}

\begin{conceptbox}
\textbf{小角三角函数}
\[
          \sum_{n=1}^{\infty}\sin\frac{\pi}{n}
        \]

因为 \(\sin(\pi/n)\sim \pi/n\)，与调和级数同敛散，所以发散。
\end{conceptbox}

\subsection*{三、根值判别法与比值判别法}

\subsubsection*{定理 9.3.3：Cauchy 根值判别法}

设 \(\sum x_n\) 是正项级数，令

\[
      r=\overline{\lim}_{n\to\infty}\sqrt[n]{x_n}.
    \]

\begin{itemize}
 \item 若 \(r<1\)，则 \(\sum x_n\) 收敛；
 \item 若 \(r>1\)，则 \(\sum x_n\) 发散；
 \item 若 \(r=1\)，则判别法失效。
 \end{itemize}

\paragraph*{证明思路}

若 \(r<1\)，取 \(q\) 使 \(r<q<1\)。充分大时 \(\sqrt[n]{x_n}<q\)，即 \(x_n<q^n\)，与几何级数比较。若 \(r>1\)，则有无穷多项满足 \(x_n>1\)，一般项不趋于 \(0\)，所以发散。

\subsubsection*{定理 9.3.4：d'Alembert 比值判别法}

设 \(\sum x_n\) 是正项级数，\(x_n\ne0\)。令

\[
      \overline r=\overline{\lim}_{n\to\infty}\frac{x_{n+1}}{x_n},
      \qquad
      \underline r=\underline{\lim}_{n\to\infty}\frac{x_{n+1}}{x_n}.
    \]

\begin{itemize}
 \item 若 \(\overline r<1\)，则 \(\sum x_n\) 收敛；
 \item 若 \(\underline r>1\)，则 \(\sum x_n\) 发散；
 \item 若 \(\overline r\ge1\) 且 \(\underline r\le1\)，则判别法失效。
 \end{itemize}

\subsubsection*{根值法与比值法的关系}

教材中的引理 9.3.1 说明，对于正项数列 \(\{x_n\}\)，有

\[
      \underline{\lim}\frac{x_{n+1}}{x_n}
      \le
      \underline{\lim}\sqrt[n]{x_n}
      \le
      \overline{\lim}\sqrt[n]{x_n}
      \le
      \overline{\lim}\frac{x_{n+1}}{x_n}.
    \]

所以：如果比值判别法能判定，根值判别法通常也能判定；但根值判别法有时能处理比值剧烈振荡的级数。

\begin{center}
\small
\begin{tabularx}{\linewidth}{|>{\raggedright\arraybackslash}p{0.18\linewidth}|>{\raggedright\arraybackslash}p{0.42\linewidth}|>{\raggedright\arraybackslash}X|}
\hline
\textbf{题型特征} & \textbf{优先方法} & \textbf{原因} \\ \hline
一般项含 \(a^n,\ n^n,\) 或整体 \(n\) 次方 & 根值判别法 & 取 \(\sqrt[n]{x_n}\) 后结构会明显简化。 \\ \hline
一般项含 \(n!, (2n)!,\) 连乘积 & 比值判别法 & \(x_{n+1}/x_n\) 会大量约分。 \\ \hline
比值极限等于 \(1\) & Raabe 或积分判别 & 根值、比值都常失效，需要更细的比较。 \\ \hline
一般项像 \(1/(n\ln^q n)\) & 积分判别 & 直接对应 \(\int dx/(x\ln^q x)\)。 \\ \hline
\end{tabularx}
\end{center}

\subsection*{四、Raabe 判别法}

当

\[
      \lim_{n\to\infty}\frac{x_{n+1}}{x_n}=1
    \]

时，根值判别法与比值判别法往往都失效。Raabe 判别法专门处理“比值非常接近 1”的正项级数。

\subsubsection*{定理 9.3.5：Raabe 判别法}

设 \(\sum x_n\) 是正项级数，且 \(x_n\ne0\)。若

\[
      \lim_{n\to\infty} n\left(\frac{x_n}{x_{n+1}}-1\right)=r,
    \]

则：

\begin{itemize}
 \item 当 \(r>1\) 时，\(\sum x_n\) 收敛；
 \item 当 \(r<1\) 时，\(\sum x_n\) 发散；
 \item 当 \(r=1\) 时，Raabe 判别法失效。
 \end{itemize}

\paragraph*{证明思路}

若 \(r>1\)，可以找到 \(t>1\)，使充分大时

\[
      \frac{x_n}{x_{n+1}}>1+\frac{t}{n}.
    \]

这会推出 \(n^t x_n\) 最终单调下降且有界，从而 \(x_n\le C/n^t\)，与 \(p\)-级数比较收敛。若 \(r<1\)，则可推出 \(n x_n\) 最终有正下界，即 \(x_n\ge c/n\)，与调和级数比较发散。

\begin{notebox}
Raabe 的记忆方式：算出来的 \(r\) 和 \(p\)-级数指数的位置很像，\(r>1\) 收敛，\(r<1\) 发散。
\end{notebox}

\subsubsection*{典型例子}

判断

\[
      1+\sum_{n=1}^{\infty}\frac{(2n-1)!!}{(2n)!!}\cdot\frac{1}{2n+1}
    \]

的敛散性。设一般项为 \(x_n\)，直接比值有

\[
      \lim_{n\to\infty}\frac{x_{n+1}}{x_n}=1,
    \]

比值法失效。但

\[
      \lim_{n\to\infty} n\left(\frac{x_n}{x_{n+1}}-1\right)=\frac32>1,
    \]

所以由 Raabe 判别法可知级数收敛。

\subsection*{五、积分判别法}

如果正项级数的一般项来自一个非负函数的取值，就可以把级数与反常积分互相转换。

\subsubsection*{定理 9.3.6：积分判别法}

设 \(f(x)\) 定义在 \([a,+\infty)\) 上，\(f(x)\ge0\)，且在任意有限区间 \([a,A]\) 上 Riemann 可积。取一列

\[
      a=a_1<a_2<\cdots<a_n<\cdots,\qquad
      u_n=\int_{a_n}^{a_{n+1}}f(x)\,dx.
    \]

则反常积分与正项级数同敛散：

\[
      \int_a^{+\infty}f(x)\,dx
      =
      \sum_{n=1}^{\infty}u_n
      =
      \sum_{n=1}^{\infty}\int_{a_n}^{a_{n+1}}f(x)\,dx.
    \]

特别地，若 \(f(x)\) 单调减少，取 \(a_n=n\)，则

\[
      f(n+1)\le
      \int_n^{n+1}f(x)\,dx
      \le f(n),
    \]

所以 \(\sum f(n)\) 与 \(\int f(x)\,dx\) 同敛散。

\subsubsection*{两个标准结论}

\begin{conceptbox}
\textbf{\(p\)-级数}
\[
          \sum_{n=1}^{\infty}\frac1{n^p}
        \]

当 \(p>1\) 收敛，当 \(p\le1\) 发散。证明来自 \(\int_1^{+\infty}x^{-p}\,dx\)。
\end{conceptbox}

\begin{conceptbox}
\textbf{对数修正级数}
\[
          \sum_{n=2}^{\infty}\frac1{n\ln^q n}
        \]

当 \(q>1\) 收敛，当 \(q\le1\) 发散。证明来自 \(\int_2^{+\infty}\frac{dx}{x\ln^q x}\)。
\end{conceptbox}

\begin{notebox}
积分判别法要求 \(f(x)\ge0\)。若没有非负性，“积分按小区间拆成级数”仍可写，但敛散性不能简单反推。
\end{notebox}

\subsection*{六、考试判别路线图}

\begin{center}
\small
\begin{tabularx}{\linewidth}{|>{\raggedright\arraybackslash}p{0.18\linewidth}|>{\raggedright\arraybackslash}p{0.42\linewidth}|>{\raggedright\arraybackslash}X|}
\hline
\textbf{你看到的形状} & \textbf{先试什么} & \textbf{常见参照} \\ \hline
一般项不趋于 \(0\) & 必要条件 & 直接发散。 \\ \hline
有理函数 \(\frac{P(n)}{Q(n)}\) & 极限比较 & \(\frac1{n^p}\)。 \\ \hline
\(\sin(1/n),1-\cos(1/n),e^{1/n}-1\) & 等价无穷小 & \(1/n,\ 1/n^2\)。 \\ \hline
阶乘、连乘积 & 比值判别 & 算 \(x_{n+1}/x_n\)。 \\ \hline
整体 \(n\) 次方、指数幂 & 根值判别 & 算 \(\sqrt[n]{x_n}\)。 \\ \hline
比值极限为 \(1\) & Raabe 判别 & 算 \(n(x_n/x_{n+1}-1)\)。 \\ \hline
\(\ln n\) 出现在分母 & 积分判别 & \(\int dx/(x\ln^q x)\)。 \\ \hline
\end{tabularx}
\end{center}

\subsection*{七、即时判断练习}

\begin{quickcheck}
判断 \(\sum_{n=1}^{\infty}\left(\frac{n}{n+1}\right)^{n^2}\)，最自然先用哪个判别法？

\tcblower
\textbf{答案：}root
\end{quickcheck}

\begin{quickcheck}
若算出 \(\lim x_{n+1}/x_n=1\)，下一步通常优先试什么？

\tcblower
\textbf{答案：}raabe
\end{quickcheck}

\begin{quickcheck}
\(\sum_{n=2}^{\infty}\frac1{n\ln n}\) 的敛散性是？

\tcblower
\textbf{答案：}diverge
\end{quickcheck}

\subsection*{八、课后题训练}

下面按教材第 24-25 页习题整理。题面完整保留，提示只给入口；写作业时仍建议先独立算一遍，再展开提示核对方向。

\begin{exercisebox}
\textbf{习题 1 \badge{正项级数敛散性}}

讨论下列正项级数的敛散性：

\[
      \begin{aligned}
      &(1)\ \sum_{n=1}^{\infty}\frac{4n}{n^4+1};\qquad
      (2)\ \sum_{n=1}^{\infty}\frac{2n^2}{n^3+3n};\\
      &(3)\ \sum_{n=2}^{\infty}\frac1{\ln^2 n};\qquad
      (4)\ \sum_{n=1}^{\infty}\frac1{n!};\\
      &(5)\ \sum_{n=1}^{\infty}\frac{\ln n}{n^2};\qquad
      (6)\ \sum_{n=1}^{\infty}\left(1-\cos\frac{\pi}{n}\right);\\
      &(7)\ \sum_{n=1}^{\infty}\frac1{\sqrt[n]{n}};\qquad
      (8)\ \sum_{n=1}^{\infty}\left(\sqrt[n]{n}-1\right);\\
      &(9)\ \sum_{n=1}^{\infty}\frac{n^2}{2^n};\qquad
      (10)\ \sum_{n=1}^{\infty}\frac{\left[2+(-1)^n\right]^n}{2^{2n+1}};\\
      &(11)\ \sum_{n=1}^{\infty}n^2e^{-n};\qquad
      (12)\ \sum_{n=1}^{\infty}\frac{2^n n!}{n^n};\\
      &(13)\ \sum_{n=1}^{\infty}\left(\sqrt{n^2+1}-\sqrt{n^2-1}\right);\\
      &(14)\ \sum_{n=1}^{\infty}\left(2n-\sqrt{n^2+1}-\sqrt{n^2-1}\right);\\
      &(15)\ \sum_{n=2}^{\infty}\ln\frac{n^2+1}{n^2-1};\qquad
      (16)\ \sum_{n=3}^{\infty}\left(-\ln\cos\frac{\pi}{n}\right);\\
      &(17)\ \sum_{n=1}^{\infty}\frac{a^n}{(1+a)(1+a^2)\cdots(1+a^n)}\quad (a>0).
      \end{aligned}
      \]

\begin{hintbox}
(1)(2)(3)(5)(13)(14)(15)(16) 多用等价无穷小或极限比较；(4)(9)(11)(12)(17) 可优先试比值或根值；(6) 用 \(1-\cos t\sim t^2/2\)；(7) 一般项不趋于 \(0\)；(8) 用 \(\sqrt[n]{n}-1=e^{(\ln n)/n}-1\sim(\ln n)/n\)。
\end{hintbox}
\end{exercisebox}

\begin{exercisebox}
\textbf{习题 2 \badge{必要条件}}

利用级数收敛的必要条件，证明：

\[
      (1)\ \lim_{n\to\infty}\frac{n^n}{\{(n-1)!\}^2}=0;\qquad
      (2)\ \lim_{n\to\infty}\frac{(2n)!}{2^n n^{2n+1}}=0.
      \]

\begin{hintbox}
把待证数列看作某个正项级数的一般项。若能证明对应级数收敛，则一般项必须趋于 \(0\)。
\end{hintbox}
\end{exercisebox}

\begin{exercisebox}
\textbf{习题 3 \badge{Raabe 判别法}}

利用 Raabe 判别法判断下列级数的敛散性：

\[
      (1)\ \sum_{n=1}^{\infty}
      \frac{n!}{(a+1)(a+2)\cdots(a+n)}\quad(a>0);
      \]\[
      (2)\ \sum_{n=1}^{\infty}\frac1{3^{\ln n}};
      \qquad
      (3)\ \sum_{n=1}^{\infty}\left(\frac12\right)^{1+\frac12+\cdots+\frac1n}.
      \]

\begin{hintbox}
统一算 \(n(x_n/x_{n+1}-1)\)。第 (2) 也可化为 \(n^{-\ln3}\)；第 (3) 与调和和 \(H_n\) 有关，近似 \(n^{-\ln2}\)。
\end{hintbox}
\end{exercisebox}

\begin{exercisebox}
\textbf{习题 4 \badge{积分型正项级数}}

讨论下列级数的敛散性：

\[
      (1)\ \sum_{n=1}^{\infty}\int_0^{1/n}\sqrt{\frac{x}{1-x}}\,dx;
      \]\[
      (2)\ \sum_{n=1}^{\infty}\int_{n\pi}^{2n\pi}\frac{\sin^2 x}{x^2}\,dx;
      \]\[
      (3)\ \sum_{n=1}^{\infty}\int_0^{1/n}\ln(1+x)\,dx.
      \]

\begin{hintbox}
先估计每个积分块的数量级。例如 \(x\to0\) 时 \(\sqrt{x/(1-x)}\sim\sqrt x\)，\(\ln(1+x)\sim x\)。
\end{hintbox}
\end{exercisebox}

\begin{exercisebox}
\textbf{习题 5 \badge{Euler 常数}}

利用不等式

\[
      \frac1{n+1}<\int_n^{n+1}\frac{dx}{x}<\frac1n
      \]

证明极限

\[
      \lim_{n\to\infty}\left(1+\frac12+\frac13+\cdots+\frac1n-\ln n\right)
      \]

存在。此极限称为 Euler 常数 \(\gamma\)。

\begin{hintbox}
令 \(a_n=H_n-\ln n\)。用题给不等式证明 \(\{a_n\}\) 单调，并证明它有界。
\end{hintbox}
\end{exercisebox}

\begin{exercisebox}
\textbf{习题 6 \badge{极限比较的边界情形}}

设 \(\sum_{n=1}^{\infty}x_n\) 与 \(\sum_{n=1}^{\infty}y_n\) 是两个正项级数。若

\[
      \lim_{n\to\infty}\frac{x_n}{y_n}=0
      \quad\text{或}\quad
      +\infty,
      \]

请问这两个级数的敛散性关系如何？

\begin{hintbox}
若 \(x_n/y_n\to0\)，说明 \(x_n\) 最终比 \(y_n\) 小很多；若 \(x_n/y_n\to+\infty\)，说明 \(x_n\) 最终比 \(y_n\) 大很多。只保留能由比较判别法推出的方向。
\end{hintbox}
\end{exercisebox}

\begin{exercisebox}
\textbf{习题 7 \badge{平方级数}}

设正项级数 \(\sum_{n=1}^{\infty}x_n\) 收敛，则 \(\sum_{n=1}^{\infty}x_n^2\) 也收敛；反之如何？

\begin{hintbox}
收敛的正项级数必有 \(x_n\to0\)，所以充分大时 \(0\le x_n<1\)，从而 \(x_n^2\le x_n\)。反例可从调和级数附近找。
\end{hintbox}
\end{exercisebox}

\begin{exercisebox}
\textbf{习题 8 \badge{Cauchy-Schwarz 估计}}

设正项级数 \(\sum_{n=1}^{\infty}x_n\) 收敛，则当 \(p>\frac12\) 时，级数

\[
      \sum_{n=1}^{\infty}\frac{\sqrt{x_n}}{n^p}
      \]

收敛；又问当 \(0<p\le\frac12\) 时，结论是否仍然成立？

\begin{hintbox}
用 Cauchy-Schwarz 不等式： \[
        \sum \frac{\sqrt{x_n}}{n^p}
        \le
        \left(\sum x_n\right)^{1/2}
        \left(\sum \frac1{n^{2p}}\right)^{1/2}.
        \] 当 \(0<p\le1/2\) 时，可构造反例。
\end{hintbox}
\end{exercisebox}

\begin{exercisebox}
\textbf{习题 9 \badge{单调函数与级数}}

设 \(f(x)\) 在 \([1,+\infty)\) 上单调增加，且 \(\lim_{x\to+\infty}f(x)=A\)。

\begin{enumerate}
 \item 证明级数 \(\sum_{n=1}^{\infty}[f(n+1)-f(n)]\) 收敛，并求其和；
 \item 进一步设 \(f(x)\) 在 \([1,+\infty)\) 上二阶可导，且 \(f''(x)<0\)，证明级数 \(\sum_{n=1}^{\infty}f'(n)\) 收敛。
 \end{enumerate}

\begin{hintbox}
第 (1) 是裂项；第 (2) 中 \(f''<0\) 说明 \(f'\) 单调减少，可尝试积分判别或与 \(f(n+1)-f(n)\) 比较。
\end{hintbox}
\end{exercisebox}

\begin{exercisebox}
\textbf{习题 10 \badge{参数积分生成级数}}

设

\[
      a_n=\int_0^{\pi/4}\tan^n x\,dx,\qquad n=1,2,\ldots.
      \]

\begin{enumerate}
 \item 求级数 \(\sum_{n=1}^{\infty}\frac{a_n+a_{n+2}}{n}\) 的和；
 \item 设 \(\lambda>0\)，证明级数 \(\sum_{n=1}^{\infty}\frac{a_n}{n^\lambda}\) 收敛。
 \end{enumerate}

\begin{hintbox}
用 \(a_n+a_{n+2}=\int_0^{\pi/4}\tan^n x(1+\tan^2x)\,dx\)，再令 \(t=\tan x\)。第 (2) 可估计 \(a_n\) 的衰减，或把和与积分交换后比较。
\end{hintbox}
\end{exercisebox}

\begin{exercisebox}
\textbf{习题 11 \badge{发散判据}}

设 \(x_n>0\)，且

\[
      \frac{x_{n+1}}{x_n}>1-\frac1n,\qquad n=1,2,\ldots.
      \]

证明 \(\sum_{n=1}^{\infty}x_n\) 发散。

\begin{hintbox}
这个条件说明 \(x_n\) 下降得不够快。尝试证明 \(n x_n\) 有正下界，或与 Raabe 判别法的发散部分比较。
\end{hintbox}
\end{exercisebox}

\begin{exercisebox}
\textbf{习题 12 \badge{构造更小的发散级数}}

设正项级数 \(\sum_{n=1}^{\infty}x_n\) 发散，且 \(x_n>0\)。证明：必存在发散的正项级数 \(\sum_{n=1}^{\infty}y_n\)，使得

\[
      \lim_{n\to\infty}\frac{y_n}{x_n}=0.
      \]

\begin{hintbox}
可以按部分和大小分块，把 \(x_n\) 乘上一个慢慢趋于 \(0\) 的系数。关键是让“每一块贡献仍然足够大”。
\end{hintbox}
\end{exercisebox}

\begin{exercisebox}
\textbf{习题 13 \badge{部分和技巧}}

设正项级数 \(\sum_{n=1}^{\infty}x_n\) 发散，且

\[
      S_n=\sum_{k=1}^{n}x_k.
      \]

证明级数

\[
      \sum_{n=1}^{\infty}\frac{x_n}{S_n^2}
      \]

收敛。

\begin{hintbox}
利用 \(S_n-S_{n-1}=x_n\)。常见比较是 \[
        \frac{x_n}{S_n^2}
        \le
        \frac1{S_{n-1}}-\frac1{S_n}
        \] 的裂项型估计。
\end{hintbox}
\end{exercisebox}

\begin{exercisebox}
\textbf{习题 14 \badge{Fibonacci 级数}}

设 \(\{a_n\}\) 为 Fibonacci 数列，证明级数

\[
      \sum_{n=1}^{\infty}\frac{a_n}{2^n}
      \]

收敛，并求其和。

\begin{hintbox}
若取 \(a_1=a_2=1,\ a_{n+2}=a_{n+1}+a_n\)，可设 \[
        S=\sum_{n=1}^{\infty}\frac{a_n}{2^n},
        \] 再利用递推关系或生成函数。标准结果是 \(S=2\)。
\end{hintbox}
\end{exercisebox}

来源参考：陈纪修、于崇华、金路《数学分析》第三版下册，第九章 §3。建议学习时把本节和 §1 数项级数的收敛性、§2 上极限与下极限 连起来看：正项级数的所有判别法，最终都服务于“部分和是否有界”。

如果某个判别法卡住，可以直接问我“为什么这里用根值/比值/Raabe/积分判别”，我会按你当前的解题步骤批改。
